\section{Compiling the operation}
\label{sec:ch03-compiling}
\index{operation compilation|(}

Position nine is where a validated document becomes something the engine can
execute. It is also where HotChocolate's naming is at its least helpful, and
saying so early saves you from wondering whether you have found three things or
one.

The middleware is called \code{OperationResolverMiddleware}. The field it holds
the compiler in is called \code{\_operationPlanner}. The class it calls is
called \code{OperationCompiler}. Three names, one step:

\begin{minted}{csharp}
using (_diagnosticEvents.CompileOperation(context))
{
    ...
    operation = _operationPlanner.Compile(
        operationId ?? Guid.NewGuid().ToString("N"),
        documentInfo.Hash.Value,
        context.Request.OperationName,
        documentInfo.Document,
        context);

    context.SetOperation(operation);
    inFlightOperation?.TrySetResult(operation);
    ...
}
\end{minted}

I would call the step compilation, because that is what the class doing the work
is called and what the diagnostic event is called. \enquote{Resolver} in the
middleware name is the older sense of resolving \emph{which} operation to run:
a document may hold several named operations, and the request's
\code{OperationName} picks one. That is genuinely a separate job from compiling
the one it picked, and it is why the middleware carries the name it does.
Chapter~\ref{ch:inside-hotchocolate} takes the compiler apart properly. For now
the useful fact is what does not go in.

\subsection*{Variables are not baked in}

\index{variables!coercion}%
What does not go in is variable values. The compiled operation is derived from
the document, the schema and the operation name, and nothing else in that
argument list is per-request.

You can watch this. Send one document with three different variable sets and
only the first pays for compilation:

\begin{minted}[fontsize=\footnotesize]{text}
parse - validate 2.246ms compile 1.033ms coerce 2.480ms execute 8.890ms total 16.770ms
    (document cache miss, operation cache miss, 17 resolvers)
parse - validate - compile - coerce 0.038ms execute 0.617ms total 0.710ms
    (document cache hit, operation cache hit, 13 resolvers)
parse - validate - compile - coerce 0.064ms execute 1.298ms total 1.562ms
    (document cache hit, operation cache hit, 15 resolvers)
\end{minted}

Both caches hit on the second and third, because the document did not change.
Coercion still runs, at position eleven, because the values did.

\index{variables!and persisted operations}%
This is the mechanical reason behind the habit
chapter~\ref{ch:hotchocolate-quickly} asked you to build in Postman: put the
identifier in the Variables pane rather than in the query text. A query with the
value inlined is a new document every time it changes, so it misses both caches,
gets validated again and compiled again, and cannot be registered ahead of time
at all. Chapter~\ref{ch:the-wire-completely} turns that last point into
persisted operations, where the client sends only a hash. The discipline starts
paying long before that, in cache hits.

\index{resolvers!count}%
The resolver counts on those lines drift for a reason worth naming, because
somebody will rerun this and get different numbers. Seventeen, thirteen,
fifteen: the three requests asked for different products, and those products
have twelve, eight and ten reviews. The count follows the data, not the
document.

\subsection*{Leaving the pipeline in the middle}

\index{cost analysis!GraphQL-Cost header}%
While we are between validation and execution, there is a way to stop here on
purpose, and it makes the pipeline visible from outside the process. Send any
query with the header \code{GraphQL-Cost: validate}:

\begin{minted}[fontsize=\footnotesize]{json}
{ "extensions": { "operationCost": { "fieldCost": 30, "typeCost": 4 } } }
\end{minted}

HTTP 200, and no \code{data} key whatsoever. The request was parsed, validated
and costed, and then the cost analyzer at position seven answered instead of
letting it through. The timeline for that request confirms it from the inside:
\code{execute -} and zero resolvers. The sibling value \code{report} runs the
query and attaches the same numbers alongside the data.

That field cost of 30 closes a loop opened in
chapter~\ref{ch:hotchocolate-quickly}. The catalog-with-reviews query selects
three fields backed by asynchronous resolvers, \code{products}, \code{reviews}
and \code{author}, and each carries \code{@cost(weight: "10")} in the exported
SDL. The plain record properties it also selects, \code{title}, \code{rating}
and \code{displayName}, carry no weight and contribute nothing. Ten and ten and
ten.

Nothing is being enforced here: \code{MaxFieldCost} and \code{MaxTypeCost} both
default to 1000, and 30 is not close. Chapter~\ref{ch:security-hardening} turns
the limits on and argues about where to set them. What the header is good for
today is answering \enquote{what does this query cost} without running it, which
is a question worth being able to ask of a query somebody has proposed adding to
a client. Which leaves the last thing on the list: where the resolvers actually
run.
\index{operation compilation|)}
